\documentclass[track=junior,level=1]{kidscode}
%% \gl{term}{week you first meet it}{meaning}{example (optional)}
\newcommand{\gl}[4]{\par\noindent\needspace{3\baselineskip}{\bfseries #1}\hfill{\footnotesize\color{kcgray}#2}\par\nopagebreak
  \noindent #3\ifstrempty{#4}{}{\par\nopagebreak\noindent{\small\color{kcgray!80!black}Example: #4}}\par\vspace{6pt}}
\newcommand{\letter}[1]{\par\needspace{6\baselineskip}\vspace{4pt}\noindent{\Large\bfseries\color{kcaccent}#1}\par\nopagebreak\vspace{2pt}}
\begin{document}
\kctitle{Word Book}{Every new word in Little Robots, from A to Z}

\begin{grownup}[How to use the word book]
Read a word and its meaning aloud, then ask your child to show you the word in Scratch or act it
out. The gray label says which week the word first comes up (W1 = week 1).
\end{grownup}

\letter{A}
\gl{algorithm}{W1}{A plan of exact steps to get a job done.}{The steps to make a jam sandwich.}

\letter{B}
\gl{backdrop}{W7}{The picture at the back of the stage.}{A party room behind the cat.}
\gl{bug}{W3}{A mistake in code. It is a clue!}{The cat says \enquote{Bye} before \enquote{Hi}.}

\letter{C}
\gl{costume}{W4}{A picture of a sprite. Changing costumes makes it look like it moves.}{The Ballerina has four costumes.}
\gl{credit}{W7}{Saying who made something, like a picture or a sound.}{\enquote{The cat comes from the Scratch library.}}

\letter{D}
\gl{debug}{W3}{To find a bug and fix it.}{Test, look, change one thing.}
\gl{Driver and Helper}{W1}{The Driver holds the mouse. The Helper points and reads. We swap when the timer runs out.}{}

\letter{E}
\gl{event}{W6}{Something that happens, like a click or a key press. An event can start a script.}{When this sprite clicked.}

\letter{G}
\gl{green flag}{W2}{The button that starts our scripts.}{}

\letter{K}
\gl{key}{W6}{A button on the keyboard, like the space bar or an arrow.}{}

\letter{P}
\gl{pen}{W8}{A crayon for the sprite. Pen down draws. Pen up stops.}{}
\gl{practice}{W5}{Doing something again and again to get better at it.}{}
\gl{program}{W1}{Steps for a robot or a computer to do.}{}

\letter{Q}
\gl{quarter turn}{W8}{A turn of 90. Four quarter turns bring you back.}{A square has four quarter turns.}

\letter{R}
\gl{repeat}{W4}{A block that runs the blocks inside it again and again.}{Repeat 4: next costume.}

\letter{S}
\gl{script}{W2}{Blocks snapped together. They run from top to bottom.}{}
\gl{sequence}{W1}{Steps in order, one after another.}{}
\gl{showcase}{W10}{A little show where we share what we made.}{}
\gl{sprite}{W2}{A character in Scratch that can act, like the cat.}{}
\gl{stage}{W2}{The place where the sprites act.}{}
\gl{stamp}{W8}{Leaves a picture of the sprite on the stage.}{}
\gl{storyboard}{W9}{A plan in boxes. It shows the start, the middle and the end.}{}

\letter{T}
\gl{test}{W3}{Click the green flag and watch what happens.}{}

\clearpage
\section{My own words}
Draw or write a new word, and what it means to you.
\begin{center}\renewcommand{\arraystretch}{3}
\begin{tabularx}{\linewidth}{|p{4.5cm}|X|}
\hline
\textbf{Word} & \textbf{What it means (draw or write)}\\ \hline
 & \\ \hline
 & \\ \hline
 & \\ \hline
 & \\ \hline
 & \\ \hline
\end{tabularx}
\end{center}
\end{document}
